% Adapted from class/medium/combine-two.tex for the dataclass Link.
\begin{blocksection}
\question Write a function \texttt{combine\char`_two}, which takes in a linked list of integers \texttt{lnk} and a two-argument function \texttt{fn}. It returns a new linked list where every two elements of \texttt{lnk} have been combined using \texttt{fn}.\\

\begin{lstlisting}
from collections.abc import Callable
def combine_two(
    lnk: LinkedList[int], fn: Callable[[int, int], int]
) -> LinkedList[int]:
    """
    >>> from operator import add, mul
    >>> lnk1 = Link(1, Link(2, Link(3, Link(4))))
    >>> print(combine_two(lnk1, add))
    (3 7)
    >>> lnk2 = Link(2, Link(4, Link(6)))
    >>> print(combine_two(lnk2, mul))
    (8 6)
    >>> combine_two((), add)
    ()
    """
    if ______________________________________:

        return ______________________________
    elif ____________________________________:

        return ______________________________
    combined = ______________________________

    return __________________________________
\end{lstlisting}

\begin{solution}[1in]
\begin{lstlisting}
def combine_two(
    lnk: LinkedList[int], fn: Callable[[int, int], int]
) -> LinkedList[int]:
    if not isinstance(lnk, Link):
        return ()
    elif not isinstance(lnk.rest, Link):
        return Link(lnk.first)
    combined = fn(lnk.first, lnk.rest.first)
    return Link(combined, combine_two(lnk.rest.rest, fn))
\end{lstlisting}
\end{solution}

\end{blocksection}

\begin{questionmeta}
\textbf{Teaching Tips}\par\nopagebreak[4]
Suggested Time: 10 min; Difficulty: Medium
\nopagebreak[4]
\begin{itemize}
    \item \textbf{Function type:} \lstinline{Callable[[int, int], int]}
    describes two integer arguments and an integer result. Connect to HOFs.
    \item \textbf{Base cases:} Separate empty and one-element inputs.
    Use \lstinline{isinstance} to distinguish a node from \lstinline{()}.
    \item \textbf{Recursive step:} Combine two values, recurse on the
    remaining list, and construct the result. Retain an unpaired last value.
\end{itemize}
\end{questionmeta}
